\chapteropen{01}

\begin{cele}
\item policzyć, ile godzin zostaje w dobie po odjęciu snu, i zapisać to jako budżet, a nie jako życzenie;
\item rozpisać dzień roboczy w blokach, które sumują się dokładnie do 24 godzin;
\item odczytać szablony tygodnia i wskazać, skąd bierze się wolny czas i ile go jest;
\item stosować zasadę „najpierw sen, resztę dopasuj” i nazwać jej koszt;
\item zapisać własną dobę w blokach i sprawdzić ją sumą.
\end{cele}

\section{Doba ma dwadzieścia cztery godziny}

Zacznijmy od rachunku, który da się zrobić w pamięci: $24 - 12 = 12$. Jeśli śpisz dwanaście godzin na dobę, przez resztę doby masz dwanaście godzin na wszystko inne: pracę, jedzenie, dom, chleb, ruch i odpoczynek. Nie ma trzynastej godziny. Nie ma też godziny, która przepadnie bez śladu, bo jej nie zapisałeś. Cała ta książka opiera się na tym jednym odejmowaniu.

Słowo „budżet” użyto tu dosłownie. Budżet to suma, którą wolno wydać, i lista pozycji, z których się składa. Doba jest budżetem zamkniętym: każda pozycja zabiera coś z tej samej puli i nie ma nadwyżki, z której można pokryć resztę. Kiedy zapisujesz dzień jako budżet, zmienia się pytanie. Zamiast „kiedy znajdę czas na chleb?” pytasz „które z dwunastu godzin oddaję chlebowi?”. Drugie pytanie ma odpowiedź. Pierwsze zwykle kończy się poczuciem winy.

Trzeba też powiedzieć wprost, czego ta książka nie obiecuje. Dwanaście godzin snu to więcej, niż zwykle podaje się jako zalecenie dla dorosłych: \verify{dorośli zwykle potrzebują około 7--9 godzin snu}. Regularnie długi sen bywa też \verify{objawem problemu zdrowotnego, który warto skonsultować z lekarzem}, a nie tylko wyborem stylu życia. Punktem wyjścia tej książki jest dwanaście godzin jako założenie do zaplanowania, a nie zalecenie. Jeśli po kilku tygodniach budzisz się nadal zmęczony, zasypiasz w ciągu dnia albo sen jest przerywany, to informacja dla lekarza, a nie wynik lepszego planowania. Książka nie rozstrzyga, czy dwanaście godzin jest dla ciebie właściwe. Pomaga zaplanować dobę, jeśli tak już śpisz albo chcesz to sprawdzić. Ile snu naprawdę potrzebujesz, rozpatrujemy w rozdziale 2.

\begin{falle}[Pułapka: dwanaście godzin snu nie oznacza dwunastu godzin wolnego]
Przy dwunastu godzinach snu łatwo pomyśleć, że „dużo zostaje”. Zostaje dwanaście godzin, ale nie zostaje dwanaście godzin odpoczynku. Dojazd, gotowanie, mycie się i sprzątanie zjadają je, zanim cokolwiek zostanie na przyjemności. Wolny czas w tym budżecie to ta część doby, która nie jest przypisana do żadnej innej pozycji. Jak zobaczysz w sekcji 3, bywa jej piętnaście do trzydziestu minut dziennie.
\end{falle}

\section{Budżet pracy i tego, co obok niej}

Pierwsza pozycja po śnie to zwykle praca. Osiem godzin to rama etatu, a nie zmierzony czas. Jeśli dochodzi dojazd, przerwy na kawę albo czas na przygotowanie się do dnia, każda z tych rzeczy jest osobną pozycją. Pomijanie ich to najczęstszy sposób, w jaki budżet zaczyna kłamać.

Po odjęciu snu i samej pracy zostają cztery godziny:
\[ 24 - 12 - 8 = 4 \]

Rachunek nie zmienia się w weekend, ale układ doby już tak. Jeśli w sobotę wstajesz o innej porze niż w dni robocze, to sen jest innym blokiem i trzeba go zapisać osobno. Przesunięcie pobudki o dwie godziny w niedzielę przenosi dwie godziny z jednej części doby do drugiej, ale nie tworzy nowych. Zanim uznasz, że weekend „daje” czas, sprawdź, czy wstajesz tak samo, czy tylko wcześniej lub później.

Te cztery godziny mają pomieścić wszystko, czego nie ma w pracy. W szablonie pracy zdalnej, bez dojazdu, rozkładają się tak:

\begin{itemize}
\item ruch i higiena: 1 godzina (prysznic, mycie zębów, spacer albo krótki trening);
\item posiłki: 1,5 godziny (śniadanie, obiad i kolacja, razem z gotowaniem i sprzątaniem po nich);
\item dom: 0,5 godziny (pranie, zakupy, sprawy, które nie mogą czekać do soboty);
\item chleb, czas aktywny: 0,5 godziny (zaprawa, formowanie, nadzór nad wypiekiem);
\item wolne: 0,5 godziny.
\end{itemize}

Suma: $1 + 1{,}5 + 0{,}5 + 0{,}5 + 0{,}5 = 4$. Rachunek się zgadza, a to nie przypadek: cała reszta doby jest tą czwórką. Zwróć uwagę na pozycję „chleb, czas aktywny”. Liczy się tu tylko czas, w którym ręce są przy cieście. Długie fermentacje i odpoczynek ciasta nie są w tej tabeli, bo nie zajmują twojej uwagi. Wrócimy do tego w rozdziałach 8 i 9.

\section{Szablony tygodnia}

Dzień roboczy nie jest jeden. Dlatego książka ma trzy szablony, każdy sumujący się do 24 godzin: A dla pracy zdalnej, B dla biura z dojazdem i C dla weekendu. Tabele są generowane z pliku \code{data/week.tsv}, więc liczby w tekście i w tabelach nie mogą się rozjechać:

\generated{week-templates}

Zwróć uwagę na dwie wartości w wierszu „Wolne”. W szablonie A wolne wynosi \val{week.A.free} godziny, czyli 30 minut. W szablonie B wynosi \val{week.B.free} godziny, czyli 15 minut. Sen w obu szablonach to \val{week.A.sleep} godzin, bo sen zapisujemy jako pierwszy i nie przesuwamy go.

Co to znaczy w praktyce? Przy pięciu dniach roboczych w tygodniu wolne zbiera się do $5 \cdot 0{,}5 = 2{,}5$ godziny w szablonie A oraz do $5 \cdot 0{,}25 = 1{,}25$ godziny w szablonie B. To cały tygodniowy odpoczynek, który nie jest pracą, snem, jedzeniem, domem ani chlebem. Wieczór ze znajomymi trwający trzy godziny w szablonie B nie mieści się w jednym dniu. Trzeba go wziąć z innego bloku albo przenieść na weekend.

Szablon C jest dniem weekendowym. Ma \val{week.C.free} godzin wolnego, czyli tam leży większość luzu tygodnia. To także tam warto zaplanować dłuższe pieczenie, bo nie koliduje z pracą.

\begin{przyklad}[Wtorek w szablonie B, krok po kroku]
Zaczynamy od dwunastu godzin snu i na bieżąco odejmujemy kolejne pozycje od doby:
\begin{tabular}{@{}lrr@{}}\toprule
Pozycja & Godziny & Zostaje\\\midrule
Doba & & 24\\
Sen & 12 & 12\\
Praca & 8 & 4\\
Dojazd & 1 & 3\\
Ruch i higiena & 0,75 & 2,25\\
Posiłki & 1,25 & 1\\
Dom & 0,5 & 0,5\\
Chleb, czas aktywny & 0,25 & 0,25\\
Wolne & 0,25 & 0\\\bottomrule
\end{tabular}

Na końcu zostaje zero. Dzień jest szczelny: nie ma godziny bez nazwy. Jeśli ostatnia pozycja wychodziłaby ujemna, znaczyłoby to, że plan jest niewykonalny, zanim wstaniesz z łóżka.
\end{przyklad}

\section{Zasada: najpierw sen, resztę dopasuj}

\begin{zasada}[Najpierw sen]
Sen wpisujesz pierwszy, w ustalonej godzinie zaśnięcia i wstania. Nie przesuwasz go, żeby zmieścić coś innego. Wszystko inne dopasowujesz do tego, co zostało.
\end{zasada}

Dlaczego kolejność ma znaczenie? Jeśli sen wpiszesz na końcu, dostanie on to, co zostanie po pracy, dojazdach i wieczornych sprawach. Wieczór zawsze znajdzie coś, co wydaje się pilniejsze niż sen, i najłatwiej zabrać je właśnie stamtąd. Wpisany na początku sen staje się stałą, a niedobór pozostałych godzin staje się widoczny na papierze. To drugie jest dla planowania równie ważne jak pierwsze.

Weźmy zmianę, która zdarza się w prawdziwym życiu. Dojazd do biura wydłuża się o pół godziny. W szablonie A wolne wynosi \val{week.A.free} godziny, więc zjada je całe, do zera. Kolejne pół godziny nie ma już skąd wziąć: trzeba zabrać z posiłków (1,5 do 1 godziny) albo z ruchu i higieny (1 do 0,5 godziny). Nie da się dodać godziny do doby. Da się tylko przesunąć ją z jednego bloku do drugiego, i tę decyzję warto zapisać, zamiast zostawiać ją do wieczora.

\section{Koszty tej decyzji}

Warto zacząć od tego, czego nie zauważasz. Najrzadziej liczymy czas przejścia między zajęciami: wyjście z domu, zamknięcie laptopa, zmiana ubrania, dojście do kuchni. Każde z tych przejść trwa kilka minut. W dobie z czterema godzinami poza pracą i snem kilka minut przejścia pięć razy dziennie to dwadzieścia pięć minut, czyli więcej niż całe wolne w szablonie B (\val{week.B.free} godziny, 15 minut). Dlatego w metodzie z sekcji 6 przejścia mają własny blok albo są wliczone w pozycję, do której należą, a nie znikają między wierszami.

Decyzja „najpierw sen” ma koszty i trzeba je nazwać. Pierwszym jest życie towarzyskie. Dwanaście godzin snu i osiem godzin pracy zajmują dwadzieścia godzin doby. Wszystko, co nie jest pracą ani snem, mieści się w czterech godzinach, razem z dojazdem, jedzeniem i domem. Kolacja z przyjaciółmi, wyjście na koncert czy wizyta u rodziny nie są w żadnym z trzech szablonów. Są możliwe, ale zabierają z wolnego, z domu albo z posiłków. Jeśli zdarzają się częściej niż kilka razy w tygodniu, budżet pęknie, zanim zauważysz, dlaczego.

Drugim kosztem jest ruch. Pozycja „ruch i higiena” ma godzinę, a to jedna pozycja, nie dwie. Trening trwający czterdzieści pięć minut z prysznicem po nim nie mieści się w niej bez zmiany czegokolwiek innego. Można to rozwiązać na dwa sposoby: przenieść trening do innego dnia albo zrobić krótszy ruch w tygodniu i dłuższy w weekend. Oba wybory są w porządku, ale oba trzeba zapisać. Wiedza o tym, że regularny ruch \verify{sprzyja zdrowemu snu, a efekt zależy od rodzaju wysiłku i pory dnia}, nie zmienia tego, że jego czas trzeba wziąć skądś.

\begin{uwaga}[Szablony nie są zaleceniem]
Liczby w szablonach są punktem wyjścia. Twoja praca, dojazd, rodzina i miejsce zamieszkania zmienią je od razu. Szablon jest sprawdzianem spójności: jeśli twój plan nie sumuje się do 24 godzin, coś jest w nim zmyślone albo pominięte.
\end{uwaga}

\section{Metoda: zapisz swoją dobę w blokach}

Teraz pracujesz na własnej dobie. Metoda ma sześć kroków i wszystkie są arytmetyczne.

\begin{enumerate}
\item Zapisz sen jako pierwszy blok: godzinę zaśnięcia i godzinę wstania, na przykład 22:00 do 10:00. To dwanaście godzin i nie zmieniasz ich w dalszych krokach.
\item Wpisz stałe bloki: pracę, dojazd, zajęcia, które są w kalendarzu. Każdy z nich ma godziny początku i końca.
\item Resztę rozpisz na pozycje z szablonu, ale w minutach, a nie w ułamkach. Godzina to 60 minut, pół godziny to 30, trzy kwadranse to 45. Ułamki są trudne do sprawdzania w głowie; minuty nie są.
\item Zsumuj wszystkie bloki. Wynik ma być równy 24 godzinom, czyli 1440 minutom.
\item Jeśli suma przekracza 24 godziny, zabierz najpierw z pozycji „wolne”, potem z elastycznych: posiłków, ruchu, domu. Nie zabieraj z snu. Jeśli suma jest mniejsza, nazwij brakującą godzinę, zamiast zostawiać ją bez nazwy.
\item Zapisz dzień tak, jak faktycznie przebiegł, przez siedem dni. Plan i rzeczywistość to dwie różne tabele i obie są potrzebne.
\end{enumerate}

\begin{przyklad}[Ta sama doba w minutach]
Dzień roboczy w szablonie A zapisany w minutach: sen 720, praca 480, ruch i higiena 60, posiłki 90, dom 30, chleb 30, wolne 30. Suma: $720 + 480 + 60 + 90 + 30 + 30 + 30 = 1440$. To dokładnie 24 godziny. Jeśli w rzeczywistości posiłki zajęły 105 minut, a wolne ma być 15, suma nadal wynosi 1440, bo przesunęliśmy piętnaście minut z jednej pozycji do drugiej. Różnica jest widoczna i ma nazwę.
\end{przyklad}

\section{Co ten rozdział zakłada, a czego nie wie}

Rozdział nie wie, jak dobrze będzie ci się spało przy dwunastu godzinach. Nie wie, czy twój lekarz zobaczy w tym problem, ani czy twoja praca pozwala na elastyczność. Wie tylko, że doba ma dwadzieścia cztery godziny, że sen zajmuje z nich dwanaście i że pozostałe dwanaście trzeba rozdzielić świadomie. Każdy następny rozdział to rozdzielanie kolejnych pozycji i sprawdzanie, czy suma wciąż się zgadza.

\begin{cwiczenia}
\exitem{Ile godzin zostaje w dobie po odjęciu dwunastu godzin snu? Zapisz rachunek.}{$24 - 12 = 12$. Dwanaście godzin na wszystko, co nie jest snem.}
\exitem{W szablonie A wypisz wszystkie pozycje inne niż sen i zsumuj je. Czy wychodzi 12 godzin?}{Praca 8, ruch i higiena 1, posiłki 1,5, dom 0,5, chleb 0,5, wolne 0,5. Razem 12 godzin, bo $8 + 4 = 12$.}
\exitem{Ile minut wolnego jest w szablonie B, a ile w szablonie A?}{Szablon B: \val{week.B.free} godziny, czyli 15 minut ($0{,}25 \cdot 60$). Szablon A: \val{week.A.free} godziny, czyli 30 minut.}
\exitem{Dojazd w szablonie A wydłuża się o pół godziny. Co się dzieje z wolnym i co musi ustąpić, żeby suma została 24 godziny?}{Wolne spada z 0,5 do zera i znika. Żeby suma została 24, trzeba zabrać pół godziny z innej pozycji, na przykład z posiłków (1,5 do 1) albo z ruchu i higieny (1 do 0,5). Nie można dodać godziny bez zabrania jej skądinąd.}
\exitem{Ile godzin wolnego zbiera się w tygodniu z pięciu dni roboczych w szablonie B i w szablonie A?}{B: $5 \cdot 0{,}25 = 1{,}25$ godziny. A: $5 \cdot 0{,}5 = 2{,}5$ godziny.}
\exitem{Dlaczego sen wpisujesz pierwszy, a nie jako ostatni?}{Jako ostatni sen dostaje to, co zostanie po pracy i wieczornych sprawach, a te zawsze zabierają najpierw to, co elastyczne. Wpisany na początku sen staje się stałą, a niedobór reszty jest widoczny na papierze.}
\exitem{Ile minut to 0,75 godziny i do której pozycji szablonu B należy ta wartość?}{45 minut. To pozycja „ruch i higiena” w szablonie B.}
\end{cwiczenia}

\zadanie{Zapisz swoją dobę w sześciu do ośmiu blokach, sen wpisz jako pierwszy, a sumę doprowadź do 24 godzin (1440 minut). Przez siedem dni zapisuj obok każdego bloku, ile faktycznie zajął. Miara sukcesu: każdy dzień sumuje się dokładnie do 24 godzin, a na końcu tygodnia wiesz, która pozycja przekraczała plan i o ile minut.}
